\documentclass[a4paper,11pt]{article}
\title{Mein erstes \LaTeX Dokument}
\author{Marc Kurz}
\date{\today}


\begin{document}
\maketitle

\section{Einleitung}
Hallo Welt! Dies ist mein erstes LaTeX-Dokument.
Wie in Kapitel~\ref{sec:UUK2} beschrieben, ist \LaTeX sehr super.

\section{Hauptteil}
Hallo Welt! Dies ist mein erstes LaTeX-Dokument.

\begin{itemize}
    \item Item 1
    \item Item 2
    \begin{itemize}
        \item Sub-Item 1
        \item Sub-item 2
        \begin{itemize}
            \item Noch eine Ebene
            \begin{itemize}
                \item Noch eine Ebene
            \end{itemize}
        \end{itemize}
    \end{itemize}
    \item Item 3
    \begin{enumerate}
        \item Nummer Aufzählungn 1
        \item Aufzählung 2
    \end{enumerate}
\end{itemize}

\begin{itemize}
    \item[+] Super
    \item[-] Nicht super
\end{itemize}

\subsection{Unterkapitel}
Lorem ipsum

\subsection{Noch ein Unterkapitel}
Lorem Ipsum

\subsubsection{Unter-Unter-Kapitel}
Test

\subsubsection{Unter-Unter-Kapitel Nr. 2}
\label{sec:UUK2}
Test

\paragraph{Paragraph}
Test

\subparagraph{Subparagraph}
Test

\section{Mathematische Formeln}
Die wohl bekannteste Formel:

\begin{equation}
    E = mc^2  
\end{equation}

Nochmals Einstein's Super-Formel:
\[
E = mc^2
\]

Ich kann die Formel auch inline im Text einbetten. Das sieht so aus: $E = mc^2$. Auch hübsch.

Nun eine komliziertere Formel:
\[
x^2 = \frac{n^2}{\frac{mc^2*\sqrt{5x}}{\int_{x}^{y}5x*\sqrt{4x}}}
\]

\[
42 = \sum_{x}^{k} x_i * k^3
\]

\section{Tabellen}

\begin{table}[h]
    \centering
    \begin{tabular}{|c||c|c|}
        \hline
        Zelle 1 & Zelle 2 & Zelle 3 \\ \hline\hline
        Zelle 4 & Zelle 5 & Zelle 6 \\ \hline
        Zelle 7 &         & Zelle 9 \\ \hline
    \end{tabular}
    \caption{Tabellen Beschreibung}
    \label{tab:placeholder}
\end{table}

\end{document}